\paragraph{Which design is the most robust?} At the user's request the goal was not the lightest design but the one
with the largest reserve against all eight limits (``weight is not that important, it is a building''). After the
clean-up seven designs meet all limits. Their smallest reserve (robustness) ranges from 0.30 (0061) to 0.02 (0157).
Two designs are almost level: 0061 (bottleneck steel) and 0169 (bottleneck package thickness). The mesh study
decides (section~\ref{sec:netz}):

\begin{itemize}
\item \textbf{0169} stays practically unchanged on a finer mesh (\SI{308}{} $\to$ \SI{319}{MPa}, $+4$~\%). Its
  reserve does not depend on numerical accuracy; its bottleneck is the package thickness, which depends on
  geometry, not on the mesh. It is therefore the \textbf{recommended design}: 24 coils, 13 stiff connectors per
  coil pair, about \SI{28}{kt} of steel.
\item \textbf{0061} has the largest reserve on the standard mesh but rises by $+58$~\% on the fine mesh
  (\SI{351}{} $\to$ \SI{556}{MPa}) and then exceeds the limit. Its bottleneck is exactly the mesh-dependent
  quantity.
\item \textbf{0093, 0008, 0014} rise by 34--46~\% on the fine mesh and end up partly just below, partly above the
  limit. They are candidates, but only after a more accurate analysis of the connector joints.
\item \textbf{0017, 0157} have little reserve and were not refined further.
\item \textbf{V3, V1} (surrogate proposals) are more robust than 0169 on the standard mesh (0.35 and 0.31), but on
  the fine mesh their stress rises by 45 and 66~\% (section~\ref{sec:ersatz1}). V3 thus drops to 0.21.
\end{itemize}

\paragraph{Narrow misses.} The five most robust inadmissible designs (0173, 0103, 0044, 0004, 0136) fail almost
all on steel stress or on the field error $B_n$. 0063 and 0136 were admissible on the coarse mesh and dropped out
after recomputation -- an example of how easily an unnoticed change of mesh flatters the result.

\paragraph{What distinguishes the variants.} Across all 198 designs:
\begin{enumerate}
\item \textbf{Number of coils:} no design with 16 coils was admissible. The coils become so long and so strongly
  curved that the cases bend too much, and the field error is larger. 20 and 24 coils split the force over more
  and shorter members.
\item \textbf{Coil shape before structure:} coil--plasma distance and coil length determine the stress more than
  number or stiffness of the connectors. The structure can compensate an unfavourable coil shape only to a limited
  extent.
\item \textbf{Package and nuclear heating} are almost purely geometric and are predicted very well by the surrogate
  ($R^2 \approx 0.98$). They can be screened reliably before any strength analysis.
\item \textbf{Stress} is predicted only roughly by the surrogate ($R^2 \approx 0.46$) because it depends on local
  peaks at the connector joints. Here the FE analysis remains indispensable.
\item \textbf{Softer connectors} (user idea ``weaken selectively'', first form) increase the stress. The idea has to
  be applied locally: soft ends, stiff core.
\end{enumerate}

\paragraph{How reliable is this?} The numbers come from a simplified model: cases as rectangular cross-sections,
connectors as beam bundles, supports as fixed points, only the magnetic force in cold operation (no cool-down
stress, no dead weight, no fatigue). The finer mesh shows that the peaks are not yet converged. The statement
``0169 is robust, 16 coils do not work'' is reliable because it holds on all meshes; the exact megapascals are not
yet.

\paragraph{Next steps.}
\begin{enumerate}
\item Compute 0169 with locally graded connectors (soft ends) and rounded joints; adapt stiffnesses iteratively by
  the stress ratio.
\item Test the conjecture that 0169 is mesh-stable because its connectors merge into almost continuous plates:
  compute 0169 once with non-overlapping connectors and once with a deliberately continuous plate.
\item Check mesh convergence at the connector joints of 0169 and 0093 with targeted local refinement.
\item Add cool-down stress (\SI{293}{K} $\to$ \SI{4}{K}) and dead weight.
\item Reuse the PicoGK model of 0169 as the starting shape for supports, cooling channels and potting.
\end{enumerate}
